\section{Висновки та подальші кроки}

У межах роботи було проведено комплексний аналіз набору даних Spotify Tracks Dataset із використанням методів інтелектуального аналізу даних та машинного навчання.

У процесі дослідження було виконано очищення та підготовку даних, проведено описовий та кореляційний аналіз аудіофіч, а також реалізовано кластеризацію музичних композицій. Результати показали наявність закономірностей між окремими характеристиками треків та їх популярністю.

Модель Random Forest продемонструвала кращі результати класифікації порівняно з Logistic Regression, що свідчить про складну структуру взаємозв'язків між ознаками музичних композицій. Аналіз важливості ознак показав, що найбільший вплив на популярність треків мають energy, danceability, loudness та valence.

Кластеризація дозволила виділити групи композицій зі схожими характеристиками, а PCA-візуалізація підтвердила наявність структурованих патернів у наборі даних.

Практична цінність роботи полягає у демонстрації можливостей застосування методів інтелектуального аналізу даних для аналізу музичних сервісів та побудови моделей прогнозування.

Подальшим розвитком роботи може бути:
\begin{enumerate}
  \item Використання більш складних ансамблевих моделей.
  \item Дослідження часових характеристик музичних трендів.
  \item Побудова рекомендаційної системи.
  \item Аналіз текстів пісень та їх впливу на популярність.
  \item Використання нейронних мереж для класифікації композицій.
\end{enumerate}

Отримані результати підтверджують ефективність використання сучасних методів аналізу даних для дослідження музичних платформ та прогнозування популярності контенту.
